% from ega4-2-fr.tex
\oldpage[IV-2]{89}

\begin{proof}
\textup{(ii)} Soit $Z$ le sous-préschéma réduit de $X$ ayant
$\overline{\{x\}}$ pour espace
sous-jacent; l'hypothèse entraîne que $\{x\}$ est ouvert \emph{dans} $Z$;
donc, pour tout $z\in Z$, le point générique $x$ de
$\operatorname{Spec}(\mathcal O_{Z,z})$ est isolé dans
$\operatorname{Spec}(\mathcal O_{Z,z})$; mais l'anneau
$A=\mathcal O_{Z,z}$ est un anneau local noethérien intègre, et l'hypothèse
entraîne qu'il existe un $f\in A$ tel que $A_f$ soit un corps; on sait
$(0,16.3.3)$ que cela entraîne $\dim(A)\leq1$. Comme
$\dim(\mathcal O_{Z,z})\leq1$ pour tout $z\in Z$, on a bien $\dim(Z)\leq1$.
\end{proof}

\begin{corollary}[5.1.11]
\label{IV.2.5.1.11-fr}
Si $X$ est un préschéma localement noethérien, toute partie fermée non vide de
$X$ contient un point fermé.
\end{corollary}

\begin{proof}
En effet, toute partie fermée de $X$ est constructible
$(0_{\mathrm{III}},9.1.1\text{ et }9.1.5)$ et il suffit d'appliquer (5.1.10).
\end{proof}

\begin{env}[5.1.12]
\label{IV.2.5.1.12-fr}
Soient $X$ un préschéma, $\mathcal F$ un $\mathcal O_X$-Module
quasi-cohérent \emph{de type fini}, $S=\operatorname{Supp}(\mathcal F)$ son
support, qui est \emph{fermé} dans $X$ $(0_{\mathrm I},5.2.2)$. Si, pour tout
$x\in X$, on considère $\operatorname{Supp}(\mathcal F_x)$ comme une partie
fermée du schéma local $\operatorname{Spec}(\mathcal O_x)$, on a, par
définition $(0,16.1.7)$,
$\dim(\mathcal F_x)=\dim(\operatorname{Supp}(\mathcal F_x))$; mais on a
\[
\operatorname{Supp}(\mathcal F_x)
=S\cap\operatorname{Spec}(\mathcal O_{X,x})
=\operatorname{Spec}(\mathcal O_{S,x})
\]
où, dans le dernier terme, $S$ désigne un sous-préschéma fermé quelconque de
$X$ ayant $S$ pour espace sous-jacent. Si l'on remarque que les composantes
irréductibles de $\operatorname{Spec}(\mathcal O_{S,x})$ sont les
intersections de $\operatorname{Spec}(\mathcal O_{X,x})$ et des composantes
irréductibles de $S$ contenant $x$, et correspondent biunivoquement à ces
dernières, on voit que
\begin{equation}
\dim(\mathcal F_x)=\dim(\mathcal O_{S,x})
\tag{5.1.12.1}\label{IV.2.5.1.12.1-fr}
\end{equation}
en vertu de (5.1.1.1); il résulte aussi de (5.1.2) que l'on a
\begin{equation}
\dim(\mathcal F_x)=\operatorname{codim}(\overline{\{x\}},S)
\tag{5.1.12.2}\label{IV.2.5.1.12.2-fr}
\end{equation}
si $X$ est localement noethérien.

On dit que $\mathcal F$ est \emph{équidimensionnel au point} $x\in X$ si
$\mathcal F_x$ est un $\mathcal O_{X,x}$-module équidimensionnel,
c'est-à-dire $(0,16.1.7)$ si $\operatorname{Supp}(\mathcal F_x)$ est
équidimensionnel en tant que partie fermée de
$\operatorname{Spec}(\mathcal O_{X,x})$; il revient au même de dire que
l'anneau $\mathcal O_{S,x}$ est équidimensionnel.

On appelle \emph{dimension de} $\mathcal F$ et l'on note $\dim(\mathcal F)$
la dimension du \emph{support} $\operatorname{Supp}(\mathcal F)$; il résulte
de (5.1.4) et (5.1.12.1) que l'on a
\begin{equation}
\dim(\mathcal F)=\sup_{x\in X}\dim(\mathcal F_x).
\tag{5.1.12.3}\label{IV.2.5.1.12.3-fr}
\end{equation}
Si $X=\operatorname{Spec}(A)$ est un schéma affine et si
$\mathcal F=\widetilde M$, où $M$ est un $A$-module de type fini, on a
$\dim(\mathcal F)=\dim(M)$ en vertu de $(0,16.1.7)$ et (5.1.4).
\end{env}

\begin{proposition}[5.1.13]
\label{IV.2.5.1.13-fr}
Soient $X$ un préschéma, $\mathcal F$ un $\mathcal O_X$-Module
quasi-cohérent de type fini, $x$ un point de $X$, $x'$ une générisation de
$x$ dans $X$; on a alors
\begin{equation}
\dim(\mathcal F_{x'})\leq\dim(\mathcal F_x).
\tag{5.1.13.1}\label{IV.2.5.1.13.1-fr}
\end{equation}
\end{proposition}

\begin{proof}
Cela résulte aussitôt de (5.1.12.1) et des définitions.
\end{proof}

\subsection{Dimension d'un préschéma algébrique}

